% ---------------------------------------------------------------------------
% Journal version for IEEE OJSP, IEEEtran journal class. Second paper of the
% series: the fitted generative model of drone ego-noise is the contribution,
% the estimators are the instrument. Preamble, macros and the reused
% paragraphs come from ../../2026-08_wrapup/src/journal.tex (do not edit that
% file). Tables: src/tables/*.tex, copied from the 2026-08 paper; figures are
% read from ../../2026-08_wrapup/figures/ through \graphicspath.
%
% Sourcing rule for this draft: every number carries a "% source:" comment
% naming the committed JSON or findings.md it was read from. Numbers that the
% nv2_* arms have not produced yet are \todo{} blocks, not placeholders that
% look like results.
% ---------------------------------------------------------------------------
\documentclass[journal]{IEEEtran}
\usepackage{newtxtext}
\usepackage{amsmath,amssymb,amsfonts}
\usepackage{graphicx}
\usepackage{booktabs}
\usepackage{multirow}
\usepackage{array}
\usepackage{microtype}
\usepackage{xcolor}
\usepackage{tikz}
\usetikzlibrary{positioning,arrows.meta,fit,calc}
\usepackage[hyphens]{url}
\usepackage[hidelinks]{hyperref}
\setlength{\emergencystretch}{1.5em}
\graphicspath{{figures/}{../figures/}{../../2026-08_wrapup/figures/}}

\newcommand{\revps}{\ensuremath{\,\mathrm{rev/s}}}
\newcommand{\todo}[1]{\textcolor{red}{[\textbf{TODO:} #1]}}

\begin{document}
\bstctlcite{journal-reference-control}

\title{A Fitted Generative Model of Drone Ego-Noise for Rotor-Speed
  Estimation}

\author{Dmitrii~Mukhutdinov and Lin~Wang%
  \thanks{The authors are with the School of Electronic Engineering and
    Computer Science, Queen Mary University of London, UK (e-mail:
    d.mukhutdinov@qmul.ac.uk; lin.wang@qmul.ac.uk).}}

\maketitle

\begin{abstract}
Estimating the four motor speeds of a quadrotor from its own noise is
limited by data: the public annotated flight recordings amount to a few
minutes from two airframes, and models trained on them recognize
trajectories they have already heard instead of reading the harmonic
comb. We replace the data with a model. We fit a probabilistic
generative model of drone ego-noise in two halves, both estimated from
real flights by exact likelihood. The first is a state-space model of
the four rotor-speed trajectories, with a slow and a resonant mode per
control axis, a per-rotor measurement process and a per-flight offset,
fitted to seven rigs; a diagonal Gaussian over a dimensionless
reparametrization of its 32 parameters samples unseen airframes. The
second is a model of the rotor-noise audio given a speed trajectory:
harmonic lines whose phase carries an integrated Ornstein--Uhlenbeck
shaft error and a per-line Lorentzian width, a per-order power profile,
a coloured floor with speed laws, and per-microphone gains, fitted by a
composite Whittle likelihood on bench recordings and then in flight.
We validate the fits with a tracker-parity gate: a frozen estimator
trained on real audio reads the rendered audio at
1.82\revps{} on the DREGON rig against a 2.19\revps{} bar and
1.62\revps{} on the Matrice 100 against 3.03. We then train
rotor-speed estimators on renders alone and on mixtures with real
audio, and report what transfers.
\end{abstract}

\begin{IEEEkeywords}
Rotor speed estimation, drone ego-noise, generative models, Whittle
likelihood, state-space models, synthetic training data, sim-to-real
transfer.
\end{IEEEkeywords}

\section{Introduction}
\label{sec:intro}

\IEEEPARstart{A}{udio} recorded on board a drone could support voice
interaction, acoustic monitoring and self-diagnostics. The main
obstacle is the ego-noise of the rotors, which buries other sources at
signal-to-noise ratios between 0 and $-30$ dB. This noise has a clear
structure: each rotor produces a comb of harmonics spaced by its
rotation rate, plus broadband aerodynamic noise. If the motor speeds
are known, the harmonics can be located, and Gulli et al.
\cite{gulli2025rotorcond} use motor-speed telemetry explicitly to
condition speech-enhancement models. Estimating the speeds from the
audio itself would remove the need for telemetry at inference time.
Flight-controller logs are not always exposed, they have to be aligned
with the audio, and they carry timing and value errors that matter at
harmonic precision (Section~\ref{sec:data}).

The obstacle to estimating the speeds is the annotation. The public
recordings with usable telemetry amount to about seven minutes of
powered flight from two airframes (Section~\ref{sec:data}), which is a
few dozen distinct speed trajectories. A model trained on that much
data can reach a low validation error and still be reading the wrong
cue. We showed previously that it is: a frequency-scaling probe, which
resamples a cruise clip and asks whether the prediction follows the
harmonics, returns a slope near zero for every regressor trained on
real audio alone, and the same models fail by an octave on an unseen
drone. The cue they use is the trajectory prior of the recordings they
saw. More annotated flight is the obvious remedy and it is expensive:
every minute needs an airframe, a flight, a synchronized array and a
telemetry log that survives the alignment.

The contribution of this paper is a fitted probabilistic generative
model of drone ego-noise that supplies the missing data. It has two
halves, and both are estimated from real recordings rather than
designed by hand. The first (Section~\ref{sec:traj}) is a continuous
state-space model of the four rotor-speed trajectories: a slow
Ornstein--Uhlenbeck mode and a damped oscillator per control axis, a
coloured per-rotor measurement process, and a per-flight offset,
fitted to seven rigs by an exact Kalman likelihood and judged against
the best fit of the synthesizer it replaces. The second
(Section~\ref{sec:noise}) is a model of the rotor-noise audio given a
speed trajectory: each order of each rotor is a line whose phase
carries the integrated shaft-speed error shared by all its orders and
an independent Wiener term that gives the line a Lorentzian width;
the line powers form a per-order profile; the broadband floor is a
Gaussian-process-shaped spectrum with its own speed law; and every
microphone has its own line and floor gains. It is fitted by a
composite Whittle likelihood on static bench recordings and then in
flight with the dynamics frozen (Section~\ref{sec:fit}). We validate
the fits against a frozen estimator trained on real audio, a spectral
proxy and the likelihood of a speed-matched oracle
(Section~\ref{sec:fitval}), and we sample rigs from the fits in three
ways of increasing difficulty (Section~\ref{sec:rigs}). The secondary
contribution is the estimator side: compact regressors with a
permutation-invariant loss, and multi-pitch salience architectures
from music transcription equipped with a per-rotor output
representation that makes them competitive (Section~\ref{sec:est}).

Section~\ref{sec:exp} is the experiment that the model exists for. We
train the estimators on renders alone, on renders mixed with the real
pool, and on a curriculum, and we score everything on the same frozen
real validation protocol against real-only references. Every training
stream in this paper is synthetic or mixed; the validation protocol
never changes.

\section{Related work}
\label{sec:related}

\subsection{Drone ego-noise}
Existing approaches to drone ego-noise
reduction include multichannel blind source separation
\cite{wang2020bss}, learned time--frequency masks for single- and
multichannel processing \cite{wang2020dnntf}, and single-channel
band-split attention \cite{liu2025edgebsrof}. These methods use the
microphone signals only. Gulli et al. \cite{gulli2025rotorcond}
condition the enhancement model on the rotor speeds taken from
telemetry. In this work we estimate those speeds from the recorded
sound, and we build the training data for that estimator.

\subsection{Pitch estimation and transcription}
Multi-pitch estimators
predict several simultaneous fundamental frequencies. Neural methods
often output a \emph{salience map}, a time--frequency array whose
values indicate candidate-pitch activity
\cite{bittner2017deepsalience,cuesta2020multif0,bittner2022basicpitch}.
Harmonic convolutions on a log-frequency axis have been used for piano
transcription \cite{wei2022hppnet} and monophonic pitch estimation
\cite{wei2022harmof0}. Our task is different, and we adapt these
architectures to four rotor-speed tracks. At cruise the rotor rates
lie within one octave (55--110\revps) and can be less than 1\revps{}
apart, so the main difficulty is to separate nearby combs.

\subsection{Order tracking}
An \emph{order} is a frequency component
indexed by a multiple of the shaft speed. Tacholess order tracking
estimates the speed from vibration or sound instead of a tachometer
\cite{peeters2019tacholess}, either by time--frequency ridge extraction
\cite{iatsenko2016ridges} or by iterative adaptive Vold--Kalman
filtering, which updates the frequency estimate and the filter
bandwidth together \cite{jiang2024iavkf}. The drone case adds four
simultaneous shafts with near-equal speeds and changes of flight regime
(idle, ramps, cruise, landing). The blind tracker of
Appendix~\ref{sec:tracker} builds on these methods.

\subsection{Synthetic training data and sim-to-real transfer}
Training on simulated data and testing on real measurements is a
standard answer to annotation scarcity, and its failure mode is
well documented. In robotics, \emph{domain randomization} trains on
simulated images whose nuisance parameters are sampled from wide
distributions, so that the real world is one draw among many
\cite{tobin2017domainrand}; randomizing the \emph{dynamics} rather
than the appearance transfers a control policy to hardware
\cite{peng2018dynrand}. The same two ideas appear in audio. Sound
event detection is trained on soundscapes synthesized from isolated
events and backgrounds, which gives strong labels that human
annotation cannot \cite{turpault2019dcase}; end-to-end diarization is
trained on simulated conversations whose statistics are fitted to
real meetings, and the fidelity of those statistics is what decides
the transfer \cite{landini2022simconv}; far-field recognition uses
simulated room impulse responses, whose mismatch against measured ones
is large enough that it is worth learning a correction
\cite{ratnarajah2021tsrir}. In rotating machinery, a numerical model
identified from measurements supplies the labelled fault data that
the real machine cannot \cite{kumar2025numerical}.

Our setting differs from all of these in one respect that shapes the
paper. The quantity we predict, the rotor speed, is the
\emph{input} to the generator, so a render is labelled exactly, with
no annotation error and no alignment; and the generator is not a
physical simulation but a probabilistic model fitted to the same
recordings we validate on. That makes the fidelity of the fit
measurable rather than assumed, and Section~\ref{sec:fitval} measures
it.

\section{Real data, splits and ground truth}
\label{sec:data}

\subsection{Task}
Given a mono or multichannel recording from a
drone-mounted array, we predict four motor speeds in revolutions per
second (rev/s) on the short-time Fourier transform (STFT) frame grid:
2048-sample windows and a 512-sample hop at 16 kHz, i.e. 31.25
frames/s. We do not assign the estimates to physical rotor identities.
For each frame, we match the four predictions to the four reference
speeds in the order that minimizes the error, both in training
(Section~\ref{sec:est}) and in evaluation, where we report the
matched mean absolute error (MAE). A speed error of 1\revps{} shifts
the fortieth harmonic by 40 Hz, so sub-rev/s accuracy matters for
harmonic-informed suppression.

\begin{table}[t]
\centering
\caption{The real data. Every recording is 8-channel unless marked
  mono; durations are the whole recording, of which the powered
  envelope (rotors spinning) is shorter. Roles: T = training pool of
  the real-data references; V, $n$ = frozen validation
  split, $n$ 8-s clips; H = held out for a future test. Telemetry: C =
  cleaned commanded speed, M = measured speed (tachometer grade), S =
  SDK log recalibrated as described in the text.}
\label{tab:data}
\scriptsize
\setlength{\tabcolsep}{2.5pt}
\begin{tabular}{@{}l>{\raggedright\arraybackslash}p{3.1cm}rcl@{}}
\toprule
Rig & Recording & Dur. & Role & Tel. \\
\midrule
\multirow{9}{*}{\parbox{1.35cm}{DREGON \cite{strauss2018dregon}\\MikroKopter,\\indoor}}
 & room 2, free flight & 82 s & T & C \\
 & room 2, hovering & 45 s & T & C \\
 & room 2, up/down & 54 s & T & C \\
 & room 2, rectangle & 44 s & T & C \\
 & room 2, spinning & 41 s & T & C \\
 & room 1, free flight, no source & 71 s & V, 8 & C, M \\
 & room 1, free flight, speech & 63 s & V, 7 & C, M \\
 & room 1, free flight, white noise & 65 s & V, 7 & C, M \\
 & room 1, two high-level source flights & 114 s & H & C, M \\
\midrule
\multirow{4}{*}{\parbox{1.35cm}{MD2 \cite{clayton2023embedded}\\Matrice 100,\\outdoor}}
 & FLY125 & 178 s & T & S \\
 & FLY124 & 122 s & V, 15 & S \\
 & FLY103 (mono) & 107 s & H & S \\
 & FLY108 (mono) & 99 s & H & S \\
\midrule
\multicolumn{2}{@{}l}{Training pool (T)} & 444 s & \multicolumn{2}{l}{$\times 8$ mics $\approx$ 1 h} \\
\multicolumn{2}{@{}l}{Validation (V): 37 clips $\times$ 8 mics} & 296 s & \multicolumn{2}{l}{74\,296 frames} \\
\multicolumn{2}{@{}l}{Speech for mixing: LibriSpeech train-clean-100} & 100 h & \multicolumn{2}{l}{} \\
\bottomrule
\end{tabular}
\end{table}

\subsection{Recordings and splits}
Table~\ref{tab:data} lists all the
recordings we use. DREGON \cite{strauss2018dregon} provides an
8-microphone array mounted under a MikroKopter quadrotor in indoor
free flight, with logged commanded speeds and, for the room-1
recordings, measured speeds from the motor controllers. The second
collection, MD2, contains outdoor flights of a DJI Matrice 100 with an
8-microphone ring and telemetry obtained through the manufacturer's
SDK \cite{clayton2023embedded}. The real training pool consists of the
five DREGON room-2 recordings and one MD2 flight (FLY125): 444 s of
audio, or about one hour of single-channel material if the eight
microphones are counted separately. It contains only a few dozen
distinct speed trajectories, since additional microphones do not
create new trajectories. That scarcity is the premise of this paper.
The validation split consists of 37 fixed 8-second clips
from three DREGON room-1 recordings (22 clips) and a different MD2
flight, FLY124 (15 clips), each with eight channels. This gives
74\,296 scored frames, of which 12.5\,\% are stopped-rotor frames,
3.7\,\% have a rotor between 0.5 and 30\revps{} (take-off and landing
transients), 13.5\,\% are ramps and 70.3\,\% are cruise. Fourteen of
the DREGON clips come from flights with a loudspeaker playing speech
or white noise in the room; the other 23 clips contain rotor noise
only. We hold out two DREGON source flights and two mono MD2 flights
and do not use them in this paper. Every number we report is a
validation number, and we select checkpoints on the same split.

Two of these recordings appear again as fitting data. The DREGON
room-2 recordings and FLY125 are the flights on which the noise model
of Section~\ref{sec:noise} is fitted, and the DREGON and MD2 telemetry
logs are two of the seven rigs of Section~\ref{sec:traj}. The
validation clips are not fitted on.

\subsection{Speech mixing}
During training we mix every 1-s chunk of rotor noise
with a LibriSpeech \cite{panayotov2015librispeech} utterance drawn from
train-clean-100 (100 h, 251 speakers) at a signal-to-noise ratio drawn
uniformly from $[-30, 0]$ dB. The noise is the reference level, with a
floor on the reference so that a silent chunk cannot be scaled into
loud speech. One sixth of the chunks form a zero-labelled ``silence''
arm, which contains speech and room noise but no rotor noise. Without
this arm the models learn that quiet means stopped, and they fail on
silence that carries sound. Each epoch therefore sees fresh mixtures.

\subsection{Ground truth}
The logs are imperfect in ways that matter at
harmonic precision. We made three decisions about them.
(i) The DREGON labels are the \emph{commanded} speed after a cleaning
step: we zero the leading constant segment before the motors respond
and median-filter single-sample spikes. The commanded track freezes at
its last value when the controller stops logging, so we do not trust it
through shutdown. Measured speeds exist only for the room-1
recordings, from a period-counting tachometer with a 49.7 Hz update
rate and a 0.27\revps{} quantization step at 80\revps{}. They agree
with the cleaned command to 0.04\,\% on average (0.11\,\% in the worst
case), so we use the command for the frozen split as well. (ii)
Compared with the audio itself, the DREGON telemetry over-reports the
speed by 0.35--0.85\,\% depending on the estimator (a
harmonic-residual fit gives 0.54\,\%, with a confidence interval of
0.45--0.62\,\%). No independent measurement decides between the
candidate corrections, so we keep the DREGON labels as logged and
carry the range as an uncertainty floor of 0.3--0.7\revps{} at cruise
under every DREGON number in this paper. (iii) The MD2 SDK log lags
the audio, and the lag grows linearly with time, which indicates a
sample-clock dilation. A fixed lag leaves a 12--16 ms RMS timing
residual, while a lag-plus-slope fit (+0.65 and +0.38 ms per second
for FLY124 and FLY125) leaves 3--5 ms. The logged speeds are also about 0.7\,\% low, so we apply a
multiplicative correction ($\times 1.0070$). The correction has to be
multiplicative because it must vanish at zero speed. We estimate both corrections from the harmonic
reconstruction residual, i.e. we fit a harmonic model to the audio at
a candidate trajectory and measure the energy it leaves, and we apply
them to the training and validation labels alike. An audio-refined
annotation obtained by minimizing that residual starting from the
telemetry has a precision floor of about 0.2\revps{}. We use it as a
fitting target where stated, not as a training or validation label.

\section{The generative model}
\label{sec:model}

% ---- filled in the "paper: section 4" commit ----

\section{Rotor-speed estimators}
\label{sec:est}

% ---- filled in the "paper: sections 5-6" commit ----

\section{Experiments}
\label{sec:exp}

% ---- filled in the "paper: sections 5-6" commit ----

\section{Discussion}
\label{sec:discussion}

% ---- filled in the "paper: discussion, conclusion, appendices" commit ----

\section{Conclusion}
\label{sec:conclusion}

% ---- filled in the "paper: discussion, conclusion, appendices" commit ----

\appendices

\section{The two-stage blind tracker}
\label{sec:tracker}

% ---- filled in the "paper: discussion, conclusion, appendices" commit ----

\section{Fit tables}
\label{sec:fittables}

% ---- filled in the "paper: discussion, conclusion, appendices" commit ----

\clearpage
% Reset column height after the final double-column float page.
\onecolumn
\twocolumn
\IEEEtriggeratref{10}
\bibliographystyle{IEEEtran}
\bibliography{references}

\end{document}
